\scanpage{67}
\mspage{9}
étale au dessus de \inserted{$\overline\Sigma^*=$} $U_{0,3}$, \deleted{revient à}\inserted{i.e.} la donnée d'un revêtement ramifié de $\Sigma$ étale au dessus de $U_{0,3}$, plus une op. de $\{1,\tau\}$ dessus compatible avec l'opération sur $\Sigma$, revient à la donnée d'une \uline{donnée de descente} de $\mathbb C$ à $\mathbb R$ dans un revêtement ramifié, étale sur $U_{0,3\,\mathbb C}$, du \deleted{$\unsure{\Sigma}$}\inserted{$\mathbb C$-}schéma $\Sigma_{\mathbb C}$. Donc la catégorie de ces revêtements mixtes \inserted{est} équivalente à celle des revêtements ramifiés du schéma \deleted{$U_{0,3}$}\inserted{$\Sigma_{\mathbb R}=\mathbb P^1_{\mathbb R}$}, étales au dessus de $\deleted{U_{0,3}}{}_{\mathbb R}=\overline\Sigma_{\mathbb R}^*=\Sigma_{\mathbb R}\smallsetminus\{0,1,\infty\}$. On trouve ainsi un \uline{foncteur pl. fidèle} canonique
\diagram{\begin{tikzcd}[ampersand replacement=\&,column sep=large]
\begin{gathered}
\text{(26)}\quad\text{cartes triangulées pondérées}\\
\inserted{\text{finies}}\\
\text{de type }\{0,1,\infty\}
\end{gathered}
\arrow[r, "\unsure{\colon}" description]
\&
\begin{gathered}
\text{revêtements }\inserted{\text{finis}}\text{ étales}\\
\text{du schéma }\mathbb P^1_{\mathbb R}\smallsetminus\{0,1,\infty\}\\
(\simeq\text{ rev. ramifiés de}\\
\mathbb P^1_{\mathbb R},\ \text{étales au dessus}\\
\text{de }U_{03\,\mathbb R}^{\unsure{\circ|*}}=\mathbb P^1_{\mathbb R}\smallsetminus\{0,1,\infty\}
\end{gathered}
\end{tikzcd}}

L'image essentielle est formée des rev. $X'$ tels que \deleted{$\unsure{X'(\mathbb R)=}$}
\[
\text{(27)}\qquad X'(\mathbb R)=\varnothing
\]
interprétation qui ne fait que traduire (25) puisque
\[
X'(\mathbb R)=\underbrace{X'(\mathbb C)}_{\substack{\text{noté }X'\\\text{dans (25)}}}^{\tau}.
\]
Le foncteur en sens inverse est donné en associant à tt rev. \inserted{ramifié} $X'$ de $\mathbf X$, étale sur $U_{0,3}$, \unsure{l'esp. top.|l'ens. top.} $X'(\mathbb C)$, muni de l'image inverse $\left.X'(\mathbb C)\right|_{(\mathbb P^1_{\mathbb R}=\overline\Sigma_{\mathbb R})}$ de l'ens. des pts réels de $\mathbb P^1(\mathbb C)=\overline\Sigma$,
\margin{66}
